\documentclass[11pt,letterpaper]{article}
\usepackage[margin=1in]{geometry}
\usepackage[T1]{fontenc}
\usepackage{amsmath,amssymb}
\usepackage{booktabs,longtable,array}
\usepackage{xcolor,listings}
\usepackage[hidelinks]{hyperref}
\hypersetup{pdftitle={SoftWorld Research Development Record}}
\lstset{basicstyle=\ttfamily\footnotesize,breaklines=true,
  columns=fullflexible,keepspaces=true,frame=single,
  showstringspaces=false}
\newcommand{\microstrain}{\ensuremath{\mu\varepsilon}}
\newcommand{\degrees}{\ensuremath{^{\circ}}}
\title{SoftWorld: Development and Preliminary Verification of a\newline
  Disc-Based Soft-Limb Simulation with Virtual Strain Sensing\newline
  \large Research Development Record --- Working Configuration}
\author{Author: [Insert researcher name]\\
  MSET graduate student, ETID\\Texas A\&M University}
\date{Documentation snapshot: September 24, 2026}

\begin{document}
\maketitle
\begin{abstract}
This report documents the development of a ROS~2/Isaac Sim platform for
commanding a compliant disc-based limb, estimating bending strain from its
measured orientations, and reconstructing its shape in RViz. The active
configuration contains six discs, five disc-to-disc intervals, and two rings
of three virtual strain gauges. Development included position and hybrid
control scripts, hanging-base geometry, damping changes, independent
pose-derived sensing, live strain plots, and migration from a second simulated
limb to a sensor-driven visualization. A verified controller-topology mismatch
caused a requested 180\degrees{} bend to produce approximately 81.818\degrees{}:
commands retained the twelve-disc controller's eleven-interval normalization.
The working loader now converts identified legacy position targets to the
five-interval configuration. An end-to-end headless test using the original
controller measured 179.9996\degrees{} at the sixth disc. These results establish
a bounded software regression result, not physical material validation or
general constant-curvature accuracy. A prior development configuration verified
a 270.0000\degrees{} torque-admittance result using accumulated inter-disc bend
measurement; the active working configuration is now intentionally limited to
180\degrees{}. Large-angle diagonal orientation errors,
sensor-model assumptions, and the absence of physical torque calibration
remain explicitly identified limitations.
\end{abstract}

\noindent\textbf{Document status.} This is a research-oriented technical record,
not a completed thesis, peer-reviewed publication, or institutionally approved
experimental protocol. It distinguishes source-code facts, recorded numerical
tests, user observations, historical requests, and proposed future experiments.
Author identity, advisor, research objectives, and thesis-specific requirements
must be supplied by the researcher. No physical experiment is claimed.

\tableofcontents
\newpage

\section{Research context, objectives, and scope}
The researcher identified this project as part of MSET research in ETID at
Texas A\&M University. The present engineering objective is to develop a
repeatable simulation workflow for compliant limb actuation and
orientation-derived shape sensing. The long-term concept is to reproduce the
configuration of a leader limb from sensor readings without providing those
leader commands to the reconstruction system.

The active workflow is:
\begin{center}
\begin{tabular}{c}
Position or torque command\\$\downarrow$\\
Physical disc-chain leader in Isaac Sim\\$\downarrow$\\
Six orientation-derived virtual strain readings\\$\downarrow$\\
Local curvature reconstruction and interpolation\\$\downarrow$\\
\texttt{/joint\_states} and robot-state publisher\\$\downarrow$\\
TF and sensor-driven limb visualization in RViz
\end{tabular}
\end{center}

The word ``physical'' above means the physics-simulated leader, not laboratory
hardware. The RViz model is a kinematic visualization, not a dynamically
actuated second robot. The earlier two-limb Isaac configuration is retained
as a backup. Earlier quadruped work is recorded as development history;
the quadruped is separate from the currently verified six-disc workflow.

\section{Evidence and documentation provenance}
This report is based on the development conversation, the repository README,
inspection of the working Python files, and available verification logs.
The supplied Webster--Jones review was consulted for the conceptual distinction
between robot-independent constant-curvature kinematics and robot-specific
actuation mappings~\cite{webster2010}. Equations labeled as implementation
equations describe this repository and should not be attributed wholesale
to that review.

The repository's base Git revision at documentation time was
\begin{lstlisting}
e3e976c97e63596887ccb22e03ad33523c25db72
\end{lstlisting}
The working tree contained additional modified and untracked files, including
the latest topology correction. Consequently, that revision alone does not
reproduce this snapshot. A durable research release should commit the intended
source, USD assets, documentation, and validation records together and record
the resulting revision and asset hashes.

The numerical excerpts in Appendix~\ref{app:results} are copied from available
test output. They are not averages across repeated trials. Historical changes
without a retained quantitative test are described as development events, not
established performance improvements.

\section{Development history and disposition}
\begin{longtable}{p{0.23\linewidth}p{0.70\linewidth}}
\toprule
Development stage & Purpose and current disposition\\\midrule
\endhead
Initial limb control & Investigated initial bending, loss of commanded pose,
and returning toward rest. Holding position, stiffness, and damping became
recurring requirements. Early intermediate values are not reconstructed here
without retained evidence.\\
Saved limb and quadruped & Requested new USD versions and reuse in a soft
quadruped. A ROS publisher failure exposed use of a dictionary for
\texttt{\_publishers}, which conflicts with rclpy's internal publisher list.
Ground hanging constraints, body curvature, and jelly-like motion were also
addressed during earlier development. The quadruped remains a separate,
historical implementation, not a validated deployment of the current limb.\\
Hanging limb and launch & Moved the fixed base into the air and hung the limb
downward. Corrected the workflow around missing installed loader paths by
rebuilding and locating scripts through the installed package.\\
Control alternatives & Separated position control from hybrid control:
phi position versus phi torque, with theta selecting the bending plane.
Several direction-range and correction attempts were rolled back. Hybrid
control remains preliminary.\\
Constant curvature and frames & Used the continuum-kinematics review to
guide bending conventions. Disc-axis visualization was requested, added,
and subsequently rolled back. Disc axes are not an active feature. The
intersection of disc radial lines was raised as a geometric accuracy issue.\\
Unexpected startup bend & Investigated possible lingering ROS controllers
after forced VS Code termination. Startup targets are zero, but live external
commands can still bend the limb; zero startup targets are not command isolation.\\
Working-file policy & Established a rule to make subsequent implementation
changes in files named \texttt{working}, preserving older files as backups.\\
Virtual sensing & Added six pose-derived virtual gauges that respond to
actual motion, including manual manipulation without a controller.\\
Follower and RViz & Initially used a second simulated limb; retained that
version as an in-simulation backup, then replaced the active follower with
sensor-driven RViz reconstruction. Added robot description, joint states,
robot-state publisher, and a map-to-root transform.\\
Visualization and plots & Addressed RViz/Snap library conflicts and the
missing map frame. Added six live strain plots and portable ZIP/build
instructions in the README.\\
Holding and accuracy & Tuned acceleration drives, body damping, solver
settings, and gravity compensation to improve holding in the working limb.
This does not establish physical material stiffness.\\
Six-disc transition & Reduced twelve discs to six, after rollback/reapplication,
and generated dedicated six-disc assets. An asset-generation crash motivated
native Isaac generation rather than relying on a renderer-free flattened asset.\\
Shorter spacing & Reduced a temporary 55~mm pitch to 25~mm, changing clear
spacing from 50~mm to 20~mm while retaining six discs. Updated geometry,
anchors, gains, and RViz spacing together.\\
Topology correction & The user measured 81.8181568\degrees{} for a requested
180\degrees{}. Identified the eleven-versus-five normalization mismatch;
added graph-identified legacy position conversion and verified the sixth disc.\\
\bottomrule
\end{longtable}

This chronology is thematic rather than a timestamped commit log. Rolled-back
experiments must not be presented as features of the current configuration.

\section{Simulation model and current geometry}
The limb is a finite rigid-disc/virtual-link chain with compliant joint drives.
It is not a finite-element model of a continuous elastomer and does not provide
a resolved material stress or strain field. Continuum terminology describes
its intended shape approximation and sensing model.

\begin{table}[htbp]
\centering
\begin{tabular}{lll}
\toprule
Quantity & Symbol & Working value\\\midrule
Disc count & $N$ & 6\\
Disc-to-disc interval count & $M=N-1$ & 5\\
Disc thickness & $t$ & 0.005~m\\
Clear spacing & $g$ & 0.020~m\\
Disc-center pitch & $d=t+g$ & 0.025~m\\
Straight base-to-tip nominal length & $L_0=Md$ & 0.125~m\\
Gauge/disc radius & $r$ & 0.020~m\\
Hanging generator base elevation & $z_0$ & 0.400~m\\
Runtime disc mass constant & $m_d$ & 0.010~kg\\
Virtual-link mass constant & $m_v$ & 0.010~kg\\
Virtual-link separation & $h$ & 0.0005~m\\
\bottomrule
\end{tabular}
\caption{Source configuration, not measured hardware dimensions or calibrated
material properties. Mass constants are also used in gravity compensation;
asset mass agreement should be checked for each research release.}
\end{table}

The hanging generator places disc $i$ at
$\mathbf{p}_i=(0,0,z_0-id)$ in the straight configuration. The base is kinematic;
the subsequent discs are movable. Numbering is zero-based: the first disc is
\texttt{Disk\_00}, and the sixth/final disc is \texttt{Disk\_05}, not
\texttt{Disk\_06}. The Isaac final disc is red; the current RViz final disc is
gray. Color does not identify an independent measurement channel.

The first interval has one base D6 joint. Each of the other four intervals
has two serial D6 joints, A and B, connected through a virtual link. There are
therefore nine active bending joints, not five individual D6 joints.
Translation and local rotZ are locked; rotX and rotY are driven. The working
topology definition and stage guard require exactly the six intended discs.
The rigid-disc view uses explicitly ordered paths, and tip sensing explicitly
selects index~5.

\section{Constant-curvature reference kinematics}
\subsection{Ideal continuous section}
For a planar constant-curvature section of arc length $L$, bending magnitude
$\Phi$, and bending-plane direction $\Theta$,
\begin{equation}
\kappa=\frac{\Phi}{L},\qquad R=\frac{1}{\kappa},\qquad
\mathbf{a}(\Theta)=(-\sin\Theta,\cos\Theta,0)^T.
\end{equation}
Angles in equations are radians. A reference convention with straight
extension along positive local Z has
\begin{align}
\mathbf{p}(s)&=\begin{bmatrix}
(1-\cos(\kappa s))\cos\Theta/\kappa\\
(1-\cos(\kappa s))\sin\Theta/\kappa\\
\sin(\kappa s)/\kappa
\end{bmatrix},\\
\mathbf{R}(s)&=R_z(\Theta)R_y(\kappa s)R_z(-\Theta).
\end{align}
The $\kappa\rightarrow0$ limit is a straight section. The hanging implementation
uses negative-Z translations, so position signs must be mapped consistently
when comparing it with this positive-Z reference convention.

Theta denotes the plane of bending, not an independently commanded physical
twist. Total phi is accumulated tangent-orientation change along the disc
chain, not the inclination of the base-to-tip chord. For an ideal planar
90\degrees{} arc, that chord is inclined
approximately 45\degrees{} from the initial tangent.

\subsection{Radial-line intersection criterion}
For a circular planar centerline, each disc plane should be normal to the local
tangent. The radial line \emph{in the bending plane} through each disc center
then passes through one common center of curvature. Arbitrary lines within
the disc plane do not satisfy this criterion. A straight limb has no finite
center of curvature, and nonuniform curvature generally has no single center.

A future quantitative check should fit a circle to measured disc centers,
test disc-normal/tangent agreement, and report distances from that fitted
center to the corresponding radial lines. Visual agreement alone is not an
adequate numerical constant-curvature test.

\subsection{Discrete model versus an inextensible continuous arc}
The RViz URDF implements each interval through three revolute coordinates,
with the composition
\begin{equation}
T_i=R_z(\theta_i)T_z(-d/2)R_y(\beta_i)
             T_z(-d/2)R_z(-\theta_i).
\label{eq:urdf}
\end{equation}
Its transmitted coordinates are $(\theta_i,\beta_i,-\theta_i)$.
This symmetric construction provides a disc orientation and hinge-based
position, but must not be confused with an exact continuous arc of fixed
length $d$. Its disc-center chord length is
\begin{equation}
c_i=d\cos(\beta_i/2).
\end{equation}
For a uniform planar bend, its equivalent circular radius is
\begin{equation}
R_{\mathrm{hinge}}=\frac{d}{2}\cot(\beta_i/2),\qquad
\ell_{\mathrm{equivalent}}=\beta_i R_{\mathrm{hinge}},
\end{equation}
whereas an inextensible constant-curvature interval would have
\begin{equation}
R_{\mathrm{arc}}=d/\beta_i,\qquad
c_{\mathrm{arc}}=2(d/\beta_i)\sin(\beta_i/2).
\end{equation}
These agree only in the small-angle limit.
Thus, fixed nominal pitch and equal orientation increments alone do not prove
fixed arc length during bending. The virtual-link offsets and physics leader
geometry introduce further details requiring separate positional validation.

\section{Position control and holding behavior}
\subsection{Distribution over the active topology}
For total commanded bend $\Phi_c$, the working controller calculates
\begin{equation}
\beta=\frac{\Phi_c}{M}=\frac{\Phi_c}{5},\qquad
\beta_{\mathrm{base}}=\beta,\qquad
\beta_{A,j}=\beta_{B,j}=\frac{\beta}{2}\quad(j=1,\ldots,4).
\end{equation}
For 180\degrees{}, this is 36\degrees{} at the base and 18\degrees{} at each
split joint. The algebraic sum of all nine scalar bend targets is 180\degrees{}.
These A/B values remain the ROS compatibility representation. Before PhysX
initialization, the loader deactivates the eight authored serial A/B constraints
and creates one direct D6 joint between each adjacent disc pair. Both A and B
half-angle callbacks address that same direct joint; their vectors are summed,
so the physical runtime target is $\beta$ at every one of the five intervals.
The unused virtual links are retained as kinematic compatibility objects. The
two half-angle vectors are summed before writing the shared direct
drive. This removes the serial-pair high-curvature constraint while preserving
the existing controller topics.
The position controller clamps total requests to $\pm180$\degrees{}. Runtime
X/Y D6 limits are $\pm80$\degrees{} so the torque-admittance mode can exceed
180\degrees{} without clipping the base interval.

Each drive receives the implementation mapping
\begin{equation}
q_x=-\beta_j\sin\Theta,\qquad q_y=\beta_j\cos\Theta.
\label{eq:drive}
\end{equation}
This is not generally equivalent to one exact axis-angle rotation at finite
angles. Physical joint coordinate coupling matters for diagonal bends;
Section~\ref{sec:validation} reports the observed limitation.

\subsection{Drive and solver configuration}
\begin{center}
\begin{tabular}{ll}
\toprule
Parameter & Working runtime value\\\midrule
Drive type & Acceleration (mass-normalized)\\
All five direct-joint stiffness / damping & 6400 / 240\\
Linear / angular body damping & 2.0 / 10.0\\
Physics frequency & 240~Hz\\
Position / velocity solver iterations & 32 / 8\\
Controller publishing frequency & 20~Hz\\
Compatibility half-target slew limit & 15\degrees{}/s\\
Effective direct-interval slew limit & 30\degrees{}/s\\
Torque magnitude limit & 1.5~N\,m (software bound)\\
Torque slew limit & 2.0~N\,m/s (software bound)\\
\bottomrule
\end{tabular}
\end{center}
These are simulator drive settings, not measured Young's modulus, beam
rigidity, or a calibrated physical damper. The torque bounds are not
hardware-safe operating specifications. The earlier $10^6$~N\,m software
bound and 1000~N\,m example were removed because they drove the model directly
to its angular stops and concealed changes in commanded torque.

In position mode, the loader applies upward forces equal to the configured
disc and virtual-link weights, using $g=9.81$~m/s$^2$. This explicitly
compensates gravity; the reported holding test should therefore not be
described as uncompensated, payload-bearing behavior. Compensation is enabled
in both position and calibrated torque-admittance modes. Consequently, neither
mode is an uncompensated hanging-load experiment.

Targets are ramped locally rather than corrected instantaneously. Since base
and split targets have different magnitudes but share the same local slew
rate, uniform final targets do not guarantee uniform curvature during the
transition. No settling-time sweep or transient constant-curvature proof
has yet been recorded.

\section{Confirmed twelve-to-six-disc normalization failure}
The preserved original position controller has eleven intervals. If it
publishes a 180\degrees{} total request to a five-interval leader, the final
disc receives only
\begin{equation}
\Phi_{\mathrm{actual}}\simeq5\left(\frac{180\degrees}{11}\right)
  =81.81818\degrees.
\end{equation}
The user's reported tip reading was 81.81815680578234\degrees{}, numerically
consistent with this mismatch. The result is a command-normalization problem,
not evidence that additional stiffness is required.

The working loader now identifies the known position controller through the
ROS publisher graph. A complete contiguous set of A/B publisher pairs is used
to distinguish five-interval from eleven-interval sources. For an identified
legacy source,
\begin{equation}
q_{\mathrm{adapted}}=q_{\mathrm{received}}\frac{11}{5}.
\end{equation}
Scaling happens before the local position clamp. A warning reports conversion.
The working five-interval controller is not scaled again, and torque targets
are not adapted. Manual publishers with a different identity are unchanged.

The local source-built Humble rclpy delivers subscription messages without
publisher-message metadata. The adapter therefore requires an unambiguous
publisher for each topic and verifies node identity through graph endpoints.
Incomplete discovery is retried. Multiple publishers or duplicate controller
identities make conversion ambiguous; a warning/error is issued and automatic
conversion is disabled. This is compatibility handling, not a guarantee against
all conflicting ROS commands. Run one controller only. It recognizes the
specified original topology and controller identity, not arbitrary remapped
or renamed controllers.

\section{Torque-commanded admittance control: implemented interface}
The hybrid script publishes a scalar bending-torque command $\tau_c$ and a
bending-plane direction $\Theta$. The current loader implements a calibrated
quasi-static admittance law,
\begin{equation}
\Phi_{\mathrm{eq}}=180\degrees\frac{\tau_c}{1.0~\mathrm{N\,m}},
\qquad |\tau_c|\leq1.0~\mathrm{N\,m}.
\end{equation}
The resulting equilibrium angle is distributed exactly as in position mode:
one fifth at each direct disc-to-disc runtime joint. The ROS compatibility
arrays still carry one tenth on both A and B topics, and the loader vector-sums
the two halves into the full direct-interval target. The same strong
acceleration drives then enforce the requested uniform-curvature equilibrium.
Torque command slew is limited to 2.0~N\,m/s.

The command scale was increased by a factor of 100 from the earlier
0.010~N\,m-per-180\degrees{} mapping. The current 0--1.0~N\,m range is
approximately 0--10.2~kgf\,cm and is therefore easier to compare with compact
actuator ratings. This rescaling changes only the admittance interface and not
the simulated equilibrium motion; it is not an experimentally identified
moment--curvature law or a claim that a physical limb requires this torque.

This is a \emph{torque-commanded admittance controller}, not raw open-loop
body-torque actuation and not a measured physical constitutive law. Earlier
raw-torque diagnostics produced 3.5555\degrees{} at 0.001~N\,m and
7.1150\degrees{} at 0.002~N\,m, but curvature was nonuniform and depended on
gravity and effective joint inertia. A temporary low-stiffness force-spring
trial also concentrated curvature toward distal intervals. Those approaches
were rejected because they did not reproduce the required constant-curvature
shape.

The active rescaled direct-topology bounded tests measured 18.0000\degrees{},
90.0000\degrees{}, and 179.9999\degrees{} at 0.1, 0.5, and 1.0~N\,m. All five
interval rotations were respectively approximately 3.6\degrees{},
18\degrees{}, and 36\degrees{}.
Negative torque produces the opposite signed target. These values verify the
implemented software mapping; they must not be interpreted as hardware torque
characterization. An explicit over-limit regression requested 1.5~N\,m; the
controller applied 1.0~N\,m and measured 179.9999\degrees{}, with
0.0001\degrees{} mapping error and approximately 36\degrees{} at every
interval. This confirms the active 180\degrees{} cap in the tested path.

The prior serial A/B topology reached only approximately 262.3\degrees{} for a
270\degrees{} request and stopped near 52\degrees{} per later disc interval.
Collision filtering, increased drive gains, and asymmetric A/B targets did not
remove this ceiling and were not retained. Replacing each runtime serial pair
with one direct disc-to-disc D6 joint produced a historical 270.0000\degrees{}
result, with all five intervals at approximately 54.0000\degrees{}. That
development result is retained as historical evidence; the active command is
clamped to 1.0~N\,m and 180\degrees{}.

While a nonzero torque command is active, incoming bend-position callbacks are
ignored, preventing a stale 20~Hz position publisher from zeroing the torque
command. The loader logs a warning when it first rejects a position command.
Publishing zero torque releases this priority. One controller at a time remains
the required experimental practice.

\section{Six virtual strain gauges}
\subsection{Placement and retained topic names}
\begin{center}
\begin{tabular}{llll}
\toprule
Station & Disc & $s/L_0$ & Angles about circumference\\\midrule
\texttt{one\_third} & \texttt{Disk\_02} & 0.40 & 0\degrees{},120\degrees{},240\degrees{}\\
\texttt{two\_thirds} & \texttt{Disk\_03} & 0.60 & 180\degrees{},300\degrees{},60\degrees{}\\
\bottomrule
\end{tabular}
\end{center}
The original placement requirement was one-third and two-thirds of the limb
length. With six discs, exact thirds lie between discs; the working rings
use the nearest symmetric pair at 40\% and 60\%. Topic names are retained
for compatibility. The second set is offset by 180\degrees{}; threefold
120\degrees{} symmetry makes the combined rings appear staggered by 60\degrees{}.

\subsection{Measured orientation to local curvature}
Let $q_i$ be the normalized disc orientation, represented as a quaternion
in $[w,x,y,z]$ order. At sensor disc $i$, the implementation forms
\begin{equation}
q_{\mathrm{rel},i}=q_{i-1}^{-1}q_{i+1},\qquad
\boldsymbol{\omega}_i=\operatorname{rotvec}(q_{\mathrm{rel},i}).
\end{equation}
Quaternion sign is chosen so that the scalar part is nonnegative; the rotation
vector consequently uses a shortest-path rotation. The local curvature
estimate is
\begin{equation}
\hat\kappa_{x,i}=\frac{\omega_{x,i}}{2d},\qquad
\hat\kappa_{y,i}=\frac{\omega_{y,i}}{2d},\qquad 2d=0.050~\mathrm{m}.
\label{eq:kappa}
\end{equation}
For the first ring this uses discs 1 and 3; for the second it uses discs 2 and 4.
Z rotation is not used in the bending-strain estimate. These are nominal-spacing
finite-difference estimates, not differentiation of a measured deformable
centerline. The orientation difference is expressed relative to the preceding
disc; interpretation as centered curvature is approximate for general
nonplanar bending.

\subsection{Curvature to virtual axial strain}
For a gauge angle $\alpha$, its nominal coordinates are
$x=r\cos\alpha$, $y=r\sin\alpha$. The implemented strain convention is
\begin{equation}
\hat\varepsilon_i(\alpha)
 =-\hat\kappa_{x,i}y+\hat\kappa_{y,i}x
 =r\left(\hat\kappa_{y,i}\cos\alpha
              -\hat\kappa_{x,i}\sin\alpha\right).
\label{eq:strain}
\end{equation}
This is dimensionless. The published reading is
\begin{equation}
\hat\varepsilon_{\microstrain}=10^6\hat\varepsilon.
\end{equation}
For example, 10,000~\microstrain{} corresponds to a strain magnitude of 0.01,
or 1\%. Positive and negative readings represent opposite bending sides
under the implementation's coordinate convention; their correspondence to
physical tensile/compressive polarity must be checked against a hardware
gauge installation and longitudinal-axis convention.

For ideal uniform curvature, substituting
$(\kappa_x,\kappa_y)=\kappa(-\sin\Theta,\cos\Theta)$ gives
\begin{equation}
\varepsilon(\alpha)=r\kappa\cos(\alpha-\Theta).
\end{equation}
The ideal three-gauge ring has zero summed bending strain. A theoretical
90\degrees{} bend over nominal length 0.125~m yields
$r\kappa\approx0.2513$, or about 251,327~\microstrain{} at the extreme gauge
position. This large value is a consequence of the assumed radius and
kinematic equation, not a calibrated prediction for a practical gauge or
proof that a real material can sustain that deformation. It is a calculated
illustration, not a retained simulation measurement.

\subsection{Independence from command and physical limitations}
Readings use live disc poses, not desired phi or theta. They therefore respond
to actual simulated motion and manually changed orientations without requiring
a controller. The gauge markers are visual objects; they do not measure a
deforming material surface.

The model excludes measured axial extension, shear strain, torsional strain,
constitutive behavior, finite-strain calibration, noise, thermal drift,
hysteresis, gauge bonding, and sensor latency. Equation~\eqref{eq:strain} is
an idealized kinematic emulation, not an experimentally validated strain field.
A material law is not inferred from the drive stiffness. Use the term
\emph{orientation-derived virtual bending strain} in research reporting.

\subsection{Virtual Wheatstone-bridge layer}
The current working implementation additionally converts each of the six
virtual strain values to an ideal independent quarter-bridge voltage. Its
explicit research assumptions are nominal resistance $R_0=350~\Omega$, gauge
factor $GF=2.0$, and excitation $V_{\mathrm{exc}}=2.5$~V. If the active gauge
occupies the selected bridge arm, its modeled resistance is
\begin{equation}
R_g=R_0(1+GF\varepsilon),
\end{equation}
and its differential output is
\begin{equation}
V_{\mathrm{out}}=-V_{\mathrm{exc}}
\frac{GF\varepsilon}{4+2GF\varepsilon}.
\label{eq:bridge}
\end{equation}
The sign is a bridge-wiring convention and must be matched to the eventual
physical connection. Near zero strain, $V_{\mathrm{out}}/V_{\mathrm{exc}}
\approx-GF\varepsilon/4$; thus 1000~\microstrain{} produces approximately
$-0.5$~mV/V, or $-1.25$~mV at 2.5~V excitation.

The voltage is published even when it is only a circuit extrapolation, while
separate validity flags are false when $|\varepsilon|>3\%$ or
$|V_{\mathrm{out}}/V_{\mathrm{exc}}|>25$~mV/V. A value with false validity
must not be reported as a feasible physical foil-gauge measurement. At
compressive strain $\varepsilon\leq-1/GF$, modeled active resistance is
nonpositive and the voltage is published as NaN. The working model's predicted
high-bend strains often exceed the 3\% flag. Consequently, this layer provides
a traceable voltage-interface hypothesis, not a validated selection of a real
gauge or a demonstration that a foil gauge can survive direct mounting.

Matplotlib normally breaks a trace at a NaN sample. For visualization only,
the working voltage plot maps these non-finite samples to a dashed
$+V_{\mathrm{exc}}/2=+1.25$~V invalid-model boundary and emits a warning. The
published ROS value remains NaN and the validity flag remains false, avoiding
the false claim that the boundary is a feasible measurement. The voltage plot
uses 0.02~V minimum vertical padding, rather than the strain plot's 1.0-unit
padding, so finite changes from 1.0 to 1.25~V remain discernible.

\section{Sensor-driven reconstruction and RViz}
Each ring supplies three strain values. Inverting Equation~\eqref{eq:strain}
uses the matrix
\begin{equation}
\mathbf{A}=r\begin{bmatrix}
-\sin\alpha_1&\cos\alpha_1\\
-\sin\alpha_2&\cos\alpha_2\\
-\sin\alpha_3&\cos\alpha_3
\end{bmatrix},\qquad
\hat{\boldsymbol{\kappa}}=\underset{\boldsymbol{\kappa}}{\arg\min}
\|\mathbf{A}\boldsymbol{\kappa}-\boldsymbol{\varepsilon}\|_2^2.
\end{equation}
Published microstrain is converted back to dimensionless strain first.
The three readings estimate two bending components; a single ring is not a
six-degree-of-freedom pose sensor. Axial stretch or torsion is not estimated.

The two recovered curvature vectors are interpolated componentwise at interval
midpoints. NumPy interpolation uses the nearest endpoint value outside the
two ring locations. Each interval is assigned
\begin{equation}
\beta_i=d\|\hat{\boldsymbol{\kappa}}_i\|,
\qquad\theta_i=\operatorname{atan2}(-\hat\kappa_{x,i},\hat\kappa_{y,i}).
\end{equation}
The generated six-disc URDF contains fifteen movable visualization joints:
three per interval as in Equation~\eqref{eq:urdf}. The reconstructed joint states
are published to \texttt{/joint\_states}; robot-state publisher converts them
to TF using the supplied robot description. A static transform connects
\texttt{map} to \texttt{soft\_limb\_disk\_00}.

The RViz limb receives sensor-derived joint states, not the leader controller's
commands. Nevertheless, since its inputs are generated from leader poses
through the same strain equation used for inversion, reconstruction agreement
is partly an internal consistency test. It is not independent validation of
the strain model. Only two rings are available; localized bends outside or
between those rings cannot generally be recovered uniquely.

\section{ROS interface and plotting}
\begin{longtable}{p{0.29\linewidth}p{0.64\linewidth}}
\toprule
Stream & Purpose and units\\\midrule
\endhead
Joint \texttt{roll/position} & Local bend target in degrees; base and pairs
01--04 A/B. Identified legacy targets are converted by the adapter.\\
Joint \texttt{yaw/position} & Bending-plane direction in degrees, wrapped
around the circumference.\\
Joint \texttt{roll/torque} & Scalar torque command in N\,m, interpreted through
the calibrated admittance law.\\
\texttt{measured\_tip/phi} & Norm of the sum of adjacent-disc local X/Y
rotation vectors, converted to degrees.\\
\texttt{measured\_tip/theta} & Direction from that accumulated vector in degrees.\\
\texttt{measured\_tip/orientation} & Relative tip quaternion $[w,x,y,z]$.\\
Ring \texttt{strain} & Three-element \texttt{Float64MultiArray}, microstrain.\\
Ring \texttt{gauge\_1..3/strain} & Individual scalar microstrain channels.\\
Ring \texttt{bridge\_voltage} & Three-element ideal quarter-bridge output in V.\\
Ring \texttt{bridge\_valid} & Three-element 1.0/0.0 nominal-range flag.\\
Ring \path{gauge_1..3/bridge_voltage} & Individual scalar voltage channels.\\
Ring \path{gauge_1..3/bridge_valid} & Individual Boolean validity channels.\\
Ring \texttt{orientation, phi, theta} & Ring pose relative to the base;
quaternion and derived angles.\\
\texttt{/robot\_description} & Generated URDF for robot-state publisher/RViz.\\
\texttt{/joint\_states}, \texttt{/tf} & Reconstructed kinematic joint states
and visualization transforms.\\
\bottomrule
\end{longtable}
The corresponding topic prefixes are:
\begin{itemize}
\item joint commands: \path{/soft_limb/pair_XX/...};
\item sensor rings: \path{/soft_limb/sensors/<station>/...}; and
\item measured tip: \path{/soft_limb/measured_tip/...}.
\end{itemize}

The working launch starts two Matplotlib plotters: one subscribes to all six
individual strain channels and one to all six bridge-voltage channels. Each
uses two rows of three plots, a rolling 20~s history, and a 20~Hz refresh limit.
The horizontal axis uses local monotonic wall time, not a synchronized
simulation timestamp. Plot histories are not automatically exported as a
research dataset. Timestamp alignment, recording, and replay need to be added
or managed through a documented logging workflow.

The station orientation topics publish each instrumented disc's quaternion
relative to the base. During the September~24, 2026 full-launch regression,
Isaac Sim initialized but the loader then exited because this relative
quaternion was referenced before assignment in the sensor-publication path.
The assignment was restored, and the headless verification harness was
extended to execute the complete pose, strain, bridge-voltage, and station
orientation publication path so this class of launch-only error is detected.
The corrected installed loader then remained active for a 35~s headless
full-entry-point test. That test also exposed a shutdown-only race: a signal
could invalidate the ROS context during rendering before the next sensor
publication. The loop now checks ROS context validity at entry and immediately
before publishing, preventing this misleading termination traceback.
A repeated 20~s test of the rebuilt installed loader reached the simulation
loop and ended at the intentional timeout without a Python traceback. The
expanded sensor-publication verifier also passed at 0.005~N\,m, with an
expected bend of 90.0000\degrees{}, measured final-second range of
90.0000--90.0000\degrees{}, and reported tracking error of 0.0000\degrees{}.

\section{Recorded preliminary verification}\label{sec:validation}
The direct runtime topology requires special treatment during command
interpolation. The base interval receives one command, whereas each of the
other four physical intervals receives the vector sum of compatibility
commands A and B. Slewing both halves at the full 20\degrees{}/s produced a
40\degrees{}/s physical target away from the base and only 20\degrees{}/s at
the base. Although the final equilibrium remained uniform, this caused a
visibly segment-wise transient. Each compatibility half is now limited to
10\degrees{}/s, so all five physical interval targets initially advanced
together at 20\degrees{}/s. The synchronized rate was subsequently increased
to 30\degrees{}/s by using 15\degrees{}/s for each compatibility half. This
preserves the commanded constant-curvature distribution during the bend while
reducing transition time. The verification harness explicitly compares all
five X/Y drive targets after the first update.
An eight-second 180\degrees{} Isaac regression at the 30\degrees{}/s setting
reported identical first-step targets $[0.0,0.125]$ degrees on all five direct
joints, a final measured range of 179.9999--179.9999\degrees{}, and
approximately 36\degrees{} at every interval (PASS).

\subsection{Method and metrics}
The headless harness opens the hanging six-disc asset, validates topology,
applies the working drive/solver configuration, and receives commands through
the actual ROS position controller. Its default run is 12~s of simulated
time at 240~Hz (2880 physics steps); it measures the last simulated second.
It does not start the complete RViz/plotter launch. Thus it is an end-to-end
command-to-physics test, not a full graphical-launch validation.
It now exercises the complete sensor-publication method, although it still
does not validate the RViz and Matplotlib graphical processes themselves.

For position tests through 180\degrees{}, the target tip quaternion is
constructed from axis $\mathbf{a}(\Theta)$ and angle $\Phi_c$. Full orientation
error is
\begin{equation}
e_R=\left\|\operatorname{rotvec}(q_{\mathrm{target}}^{-1}
q_{\mathrm{tip}})\right\|,
\end{equation}
reported in degrees. Reported phi instead sums the local X/Y rotation vector
$\operatorname{rotvec}(q_i^{-1}q_{i+1})$ over all five disc intervals and
takes the norm of the result. Each local increment remains on its shortest
quaternion path, but their sum can represent an accumulated bend above
180\degrees{}. The raw published tip quaternion remains orientation-only and
cannot by itself distinguish, for example, 0\degrees{} from 360\degrees{}.

Harness tolerances are 2\degrees{} full orientation error and a 1\degrees{}
peak-to-peak phi range during the final second. These are software-test
tolerances, not calibrated measurement uncertainty. Isaac fast shutdown can
obscure a raised exception in the process exit status; inspect the explicit
\texttt{VERIFY RESULT: PASS/FAIL} and retain the output rather than relying
on exit status alone.

\begin{table}[htbp]
\centering
\begin{tabular}{p{0.29\linewidth}p{0.20\linewidth}p{0.27\linewidth}l}
\toprule
Case & Final phi & Full tip error & Status\\\midrule
Working controller, 180\degrees{}, theta=0 & 180.0001\degrees{} & 0.0001\degrees{} & PASS\\
Legacy controller with adapter, same target & 179.9996\degrees{} & 0.0004\degrees{} & PASS\\
Torque admittance, 0.1~N\,m & 18.0000\degrees{} & 0.0000\degrees{} mapping error & PASS\\
Torque admittance, 0.5~N\,m & 90.0000\degrees{} & 0.0000\degrees{} mapping error & PASS\\
Torque admittance, 1.0~N\,m & 179.9999\degrees{} & 0.0001\degrees{} mapping error & PASS\\
Active cap, 1.5~N\,m requested & 179.9999\degrees{} & Applied 1.0~N\,m & PASS\\
Historical pre-cap torque trial, 1.5~N\,m & 270.0000\degrees{} & 0.0000\degrees{} mapping error & PASS\\
Large diagonal, 180\degrees{}, theta=45\degrees{} & About 178.93\degrees{} & About 10.25\degrees{} & Limitation\\
User observation before compatibility fix & 81.8181568\degrees{} & Not recorded & Mismatch\\
\bottomrule
\end{tabular}
\caption{Individual recorded results. The diagonal result is retained from
the prior working drive-map test and README, not a new diagonal trial of the
latest compatibility adapter. No repeated-trial statistics are claimed.}
\end{table}

The successful tests printed a final-second range of
179.9996--179.9996\degrees{} at four decimal places. This establishes no
visible variation at the reported precision under those conditions, not
mathematically zero oscillation. Adjacent-disc orientation increments were
approximately 36\degrees{} in the single-axis test. Matching orientation
increments does not by itself validate disc-center geometry, payload
robustness, hardware behavior, or all-direction performance.

Six offline adapter regressions passed: legacy scaling, working no-double-scale,
manual publisher preservation, incomplete graph detection, duplicate controller
ambiguity, and the local Humble message-only callback. A real ROS test using
the original publisher also confirmed base 36\degrees{} and an algebraic
total of 180\degrees{} across nine active drives. The ROS package rebuilt
successfully after the correction.

\section{Third objective: three-module hanging snake}
The verified single-module design was extended into a serial hanging robot
named \texttt{soft\_3\_working}. The configuration contains three six-disc
modules (18 physical discs total), five actuated intervals per module
(15 actuated intervals total), and two fixed inter-module connectors. The
uppermost disc is kinematic, leaving the complete chain suspended under
gravity. Each module receives its own bending-torque command $\tau_{c,j}$ and
bending-plane direction $\theta_j$, for $j\in\{1,2,3\}$. The generator,
loader, controller, plotter, verifier, launch, and USD filenames use
the same shortened \texttt{soft\_3\_working} naming convention. The established
\path{/soft_snake/module_N} ROS topic namespace is retained to avoid breaking
recorded topic-based experiments. The existing active
admittance law is applied independently,
\begin{equation}
  \phi_j=180\degrees{}\frac{\tau_{c,j}}{1.0~\mathrm{N\,m}},
  \qquad |\tau_{c,j}|\leq1.0~\mathrm{N\,m},
\end{equation}
with $|\phi_j|\leq180\degrees{}$. Within each module, the total target is
distributed uniformly over its five physical intervals and all five targets
use the same slew rate. The connectors are not assigned a bend command, so
they transmit the terminal pose of one module to the base of the next.

The sensing layout is repeated locally on every module. One ring of three
gauges lies at one-third module length at $0\degrees{}$, $120\degrees{}$, and
$240\degrees{}$; a second ring lies at two-thirds length at $180\degrees{}$,
$300\degrees{}$, and $60\degrees{}$. Thus the serial robot exposes 18 strain
channels, 18 ideal quarter-bridge voltage channels, and 18 validity flags.
The per-gauge namespace is
\path{/soft_snake/module_N/sensors/<station>/gauge_G/<quantity>}.
The presentation records the 18 bridge-voltage channels rather than repeating
the strain plot. Simulator screenshots use Isaac Sim's \emph{Grey Studio}
viewport light rig, which the working loader selects programmatically before
the capture. The \texttt{soft\_3\_working} ROS launch now starts the live
18-channel bridge-voltage plotter together with Isaac Sim; plotting begins as
soon as sensor messages are available.

\subsection{Automated voltage-dataset acquisition}
A dedicated ROS~2 experiment node,
\path{scripts/collect_soft_3_working.py}, now performs repeatable command sweeps
and writes the synchronized results directly to an Excel-compatible workbook.
The node publishes the three module torque commands and three bending-plane
directions, subscribes to measured $\phi_j$ and $\theta_j$, and samples all 18
individual bridge-voltage and validity topics. It waits until every sensor is
fresh before beginning and rejects concurrent use of the ordinary
\texttt{soft\_3\_working\_controller} by default, because two command
publishers would make the independent variable ambiguous. On normal exit or
keyboard interruption, all commands are returned to zero and the samples
collected so far are saved.

The default design includes a zero-command baseline and isolates one module at
a time. For every module it evaluates
$\tau_c\in\{0.25,0.50,0.75,1.00\}$~N\,m and
$\theta\in\{-90,-45,0,45,90\}\degrees{}$, producing 61 cases. Each case is
allowed to settle for 3~s and is then recorded for 1~s at 20~Hz, for a planned
total of approximately 1,220 raw time samples in 4.1~min. Torque values,
directions, repetitions, settling duration, recording duration, sample rate,
and output path are command-line parameters. A three-repetition, 2~s recording
protocol produces approximately 7,320 samples, while the optional
\texttt{all\_same} sweep applies each command pair to all modules together.

The workbook contains five sheets: \texttt{Metadata} records the protocol and
model assumptions; \texttt{Sensor\_Map} identifies module, axial station,
gauge, and circumferential angle; \texttt{Cases} records command and mean
measured pose for each trial; \texttt{Samples} preserves synchronized raw
time-series rows; and \texttt{Sensor\_Summary} reports count, mean, sample
standard deviation, minimum, maximum, and valid fraction for every sensor and
case. A small standard-library writer,
\path{scripts/xlsx_writer_working.py}, creates the XLSX archive without adding
a pandas, openpyxl, or xlsxwriter dependency. This is a simulation data
acquisition mechanism: the stored voltages remain pose-derived predictions of
the ideal quarter-bridge model, not calibrated physical measurements.

A live smoke test collected a baseline and a 0.25~N\,m, $0\degrees{}$
all-module case. All 18 voltage and validity channels became ready, 10 raw
rows were saved, and the loaded case spanned approximately
$-0.075827922$ to $+0.096603947$~V, compared with the zero-command baseline.
The generated workbook passed ZIP/XML integrity checks and was opened and
resaved by LibreOffice Calc. This verifies the software data path and workbook
compatibility, not sensor-model accuracy.

The full default sweep was executed on October~9, 2026 and retained as
\path{research_data_working/soft_3_working_voltage_dataset_20261009_094852.xlsx}.
All 61 planned cases completed and all 1,220 planned synchronized rows were
written. Across 18 sensors this produced 21,960 validity observations, of which
14,954 (68.0965\%) were inside the nominal bridge-model range and 7,006 were
flagged invalid. There were 21,877 finite voltage cells; 83 cells reached the
non-finite model boundary and were intentionally left blank while their
validity flags were retained. The finite simulated values spanned
$-0.430364539$ to $+1.249186333$~V. The XLSX archive passed integrity testing
and LibreOffice Calc opened and resaved it. The substantial invalid fraction is
itself an experimental result: it indicates that the present pose-derived
strain/ideal-bridge assumptions exceed their stated nominal operating region
for many commanded poses and must not be treated as validated hardware output.

A 12-second headless regression commanded
$(\tau_{c,1},\tau_{c,2},\tau_{c,3})=(0.50,0.75,1.00)$~N\,m and
$(\theta_1,\theta_2,\theta_3)=(0,90,-90)\degrees{}$. Measured module bends
were 90.0047476\degrees{}, 134.9999084\degrees{}, and
179.9998960\degrees{}, respectively. Corresponding absolute tracking errors
were 0.0047476\degrees{}, 0.0000916\degrees{}, and 0.0001040\degrees{}.
The final-second ranges were 0.0001642\degrees{}, 0.0000568\degrees{}, and
0.0000713\degrees{}. Residual bends across the two fixed connectors were
0.0000672\degrees{} and 0.0000394\degrees{}. All five first-step targets were
synchronized within each module, 18 gauge publishers were detected, and the
test returned \texttt{PASS}. These results validate the software realization,
not physical torque or strain calibration.

\section{Known limitations and rejected approaches}
\begin{enumerate}
\item Equation~\eqref{eq:drive} produces a substantial full-orientation error
at the recorded large diagonal bend. Exact all-direction constant-curvature
tracking is not established, despite successful single-axis tracking.
\item An attempted runtime rotation of both D6 reference frames invalidated
physics handles in the test and was removed. That strategy is not active.
\item Several earlier corrections and disc-count changes were explicitly
rolled back. Their intended benefits must not be treated as retained results.
\item Equal final orientation increments do not prove continuous fixed-length
curvature, and local slew limits may produce nonuniform transition shapes.
\item Gravity-compensated holding does not establish loaded or unassisted
holding performance. Contact, disturbance, and mass uncertainty need testing.
\item Strain estimation and reconstruction share an idealized equation;
neither has been validated against real gauges or an independent material model.
\item ROS graph-based conversion depends on publisher identity and discovery.
Stale/duplicate controllers and nonstandard remappings remain operational risks.
\item Two sensor rings cannot uniquely recover arbitrary distributed curvature,
torsion, axial stretch, or full shape without additional assumptions.
\item The current direct-topology torque-admittance law has been verified only
for unloaded, single-axis simulation. The active rescaled interface spans 0.1
to 1.0~N\,m through a maximum of 180\degrees{}. It is not raw torque
control, a material constitutive model, or a hardware calibration; negative
commands, payloads, disturbances, and arbitrary bending directions require
separate validation. The quadruped and two-physical-limb backup were not part
of these tests.
\end{enumerate}

\section{Reproducible operation and asset handling}
\subsection{Environment}
The working environment uses Ubuntu 22.04, ROS~2 Humble, Isaac Sim~4.5.0,
an NVIDIA GPU/driver, colcon, RViz2, robot-state publisher, TF2, and Python
Matplotlib/Tk. One successful diagnostic environment reported an RTX~3080 and
driver 580.178.04; this is a test environment, not a universal prerequisite.
For each recorded experiment retain exact versions, USD hashes, timestep,
solver settings, drive values, and whether gravity compensation was enabled.

Paths below are portable examples. Replace the ROS setup path for a
source-built installation and set the actual Isaac installation location.
Run from the extracted repository root named \texttt{soft\_world}.
\begin{lstlisting}
cd /path/to/soft_world
source /opt/ros/humble/setup.bash
export ISAAC_SIM_PATH=/path/to/isaacsim
colcon build --packages-select soft_world
source install/setup.bash
ros2 launch soft_world soft_world_working.launch.py
\end{lstlisting}
In a separate terminal, source the same ROS and workspace setup, then run
\begin{lstlisting}
cd /path/to/soft_world
source /opt/ros/humble/setup.bash
source install/setup.bash
python3 scripts/control_soft_limb_working.py --ros-args \
  -p phi_position:=180.0 -p theta_position:=0.0
ros2 topic echo /soft_limb/measured_tip/phi --once
\end{lstlisting}
The echo is run in another terminal while the controller is active. Default
working position parameters are phi=90\degrees{} and theta=0\degrees{};
explicit experimental parameters are preferable.

To verify independently from the graphical launch, stop other controllers,
source the environment, and run
\begin{lstlisting}
python3 scripts/test_position_topology_working.py
"$ISAAC_SIM_PATH/python.sh" scripts/verify_bend_working.py \
  --phi 180 --theta 0
"$ISAAC_SIM_PATH/python.sh" scripts/verify_bend_working.py \
  --phi 180 --theta 0 --legacy-controller
\end{lstlisting}
Retain stdout/stderr and the explicit verification status. Run simulations
sequentially to avoid resource contention. These commands require functioning
GPU access and permission to write runtime ROS/Isaac logs.

\subsection{RViz and safe operation}
Use fixed frame \texttt{map}, add RobotModel, choose description source Topic,
and set description topic \texttt{/robot\_description}. The working launch
publishes the map-to-base transform. Its RViz/plotter processes remove selected
Snap-injected GTK/GIO paths to avoid the observed host-library ABI conflict;
this workaround is environment-specific.

Start only one controller after the simulation loads. Stop the controller
before stopping the launch. A forced VS Code exit may leave controllers alive;
inspect nodes and process identity before terminating a confirmed stale
process. Restart the simulator after loader changes; a running Python process
does not automatically reload an edited installed script.

\subsection{USD generation, portability, and version control}
The active launch loads \path{usd/soft_limb_hanging_6_disc_working.usd} directly.
It does not prune a twelve-disc scene at runtime. A matching six-disc position
asset is retained separately from older twelve-disc assets.
Native generators are invoked by
\path{scripts/create_six_disc_assets_working.py}; they require Isaac's
initialized schema/runtime environment. The authored geometry updater preserves
existing native metadata rather than rebuilding a flattened asset without
that environment. Asset references must be inspected when distributing to
another machine; portable Python paths do not alone establish portable USD
references.

For ZIP sharing, retain source, launch files, package metadata, USD assets,
README, and this research documentation; omit generated \texttt{build/},
\texttt{install/}, \texttt{log/}, and caches. Rebuild at the recipient's location.
\texttt{.gitignore} only governs untracked files in Git; it does not delete
files, untrack committed assets, or filter a manually created ZIP. Preserve
the working-file policy and avoid replacing historical backups.

\section{Recommended next research experiments}
These are proposed tests, not completed results.
\begin{enumerate}
\item Run repeated position trials across phi, theta, and both bend signs.
Report quaternion error, signed interval angles where definable, final phi
range, overshoot, settling time, and run-to-run variation.
\item Compare measured disc centers and orientations with a fitted circle
and the continuous fixed-length reference. Report radial-line intersection
residuals and distinguish hinge geometry from true continuous-arc geometry.
\item Sweep damping, drive gains, solver iterations, timestep, gravity
compensation, and applied load separately. Control these variables rather
than attributing all angle errors to stiffness.
\item Test localized manual perturbations and known nonuniform curvature.
Compare sensor-driven reconstruction against independent pose ground truth,
especially outside the sensor rings.
\item Validate virtual strain against a separate deformable model or hardware
gauges, including sensor radius, sign, axial extension, finite-strain behavior,
gauge factor, noise, calibration, and temperature sensitivity as appropriate.
\item Select the physical limb material and maximum curvature before purchasing
foil gauges. A 350~ohm linear foil gauge and a six-channel quarter-bridge DAQ
are a candidate architecture, but the simulated 25--50\% high-bend surface
strain is outside ordinary foil-gauge assumptions.
\item Introduce synchronized simulation timestamps, durable datasets, and
experiment manifests. Preserve raw data and analysis code, not only screenshots
or rolling plot displays.
\item Validate hybrid effort equilibrium and torque distribution separately;
only then consider hardware transfer or quadruped integration.
\end{enumerate}
Potential research questions include whether six-disc reconstruction error
remains acceptable relative to a twelve-disc reference, how sparse gauge rings
limit shape observability, and whether geometry-aware actuation can reduce the
large-angle diagonal error without sacrificing stable compliant motion.

\section{Conclusion}
The working platform supports a six-disc hanging leader, separate position
and calibrated torque-admittance interfaces, six pose-derived strain channels,
live graphs, and sensor-driven RViz reconstruction. It now also includes the
\texttt{soft\_3\_working} serial hanging robot with three independently
directed torque-admittance modules and 18 pose-derived gauges. The confirmed eleven-versus-five
normalization failure was corrected in the loader and reproduced in a
successful legacy-controller regression. The latest reported 180.0001\degrees{}
position result and 179.9999\degrees{} active torque-admittance endpoint are
specific software verification results; the earlier 270.0000\degrees{} result
is retained only as a pre-cap development record. Physical
strain validity, exact inextensible constant-curvature geometry, arbitrary
bending-direction accuracy, and hardware applicability remain open research
tasks rather than established conclusions.

\subsection{Research progress presentation}
An editable progress presentation was created to communicate the work completed
after the successful single-module position-control milestone. It includes
fresh captures of the current Isaac Sim environment, live six-channel strain
and bridge-voltage graphs collected while stepping the torque command through
0, 0.5, and 1.0~N\,m, quantitative constant-curvature and 180\degrees{}-cap
figures, the ROS/RViz data path, the three-module snake architecture and its
18-channel sensing/validation results, the automated torque--theta XLSX
acquisition protocol, limitations, and proposed next experiments.
The source presentation, PowerPoint export, PDF export, and captured assets are
stored under \texttt{research\_documentation}; plotted sensor values are
explicitly labeled as simulated and pose-derived rather than hardware-validated.

\begin{thebibliography}{9}
\bibitem{webster2010}
R.~J. Webster III and B.~A. Jones,
``Design and kinematic modeling of constant curvature continuum robots:
A review,'' \emph{The International Journal of Robotics Research},
vol.~29, no.~13, pp.~1661--1683, 2010.
\href{https://doi.org/10.1177/0278364910368147}{doi:10.1177/0278364910368147}.
\bibitem{implementation}
SoftWorld working repository: README, working Python implementation, and
six-disc USD assets. Inspected September~15, 2026; working snapshot extends
the base revision recorded in Section~2.
\bibitem{records}
SoftWorld development conversation and headless/ROS verification output,
September~15, 2026. Numerical excerpts retained in this document and
\path{validation_record_working.txt}; not an independently published dataset.
\end{thebibliography}

\appendix
\section{Active file map}
\begin{longtable}{p{0.53\linewidth}p{0.40\linewidth}}
\toprule
Repository-relative file & Role\\\midrule
\endhead
\path{launch/soft_world_working.launch.py} & Simulator, RViz, TF, description,
and plots.\\
\path{launch/soft_3_working.launch.py} & Three-module hanging
snake launch.\\
\path{scripts/load_soft_3_working.py} & Three-module physics,
control consumption, measurements, and 18-gauge publication.\\
\path{scripts/control_soft_3_working.py} & Three independent
phi-torque and theta command pairs.\\
\path{scripts/plot_soft_3_working.py} & Eighteen strain or bridge
voltage graphs.\\
\path{scripts/verify_soft_3_working.py} & Bounded three-module
headless regression.\\
\path{scripts/collect_soft_3_working.py} & Configurable command sweep and
18-channel Excel data acquisition.\\
\path{scripts/xlsx_writer_working.py} & Dependency-free XLSX writer.\\
\path{scripts/test_xlsx_writer_working.py} & Offline workbook-format
regression.\\
\path{scripts/soft_3_working.py} & Generator for the serial
hanging USD.\\
\path{usd/soft_3_working.usd} & Active three-module hanging asset.\\
\path{scripts/load_soft_limb_working.py} & Physics configuration, command
consumption, sensing, torque application, and reconstruction.\\
\path{scripts/control_soft_limb_working.py} & Five-interval position command.\\
\path{scripts/control_soft_limb_hybrid_working.py} & Calibrated torque-admittance
and direction command.\\
\path{soft_world/limb_configuration_working.py} & Shared six-disc topology,
tip index, bend normalization, and stage validation.\\
\path{scripts/limb_configuration_working.py} & Import bridge that avoids the
legacy \texttt{scripts/soft\_world.py} package-name collision.\\
\path{scripts/position_topology_adapter_working.py} & Identified legacy position
conversion.\\
\path{scripts/test_position_topology_working.py} & Six offline adapter regressions.\\
\path{scripts/verify_bend_working.py} & ROS-to-physics headless verification.\\
\path{scripts/plot_strain_gauges_working.py} & Six live plots.\\
\path{scripts/bridge_model_working.py} & Reproducible ideal 350~ohm
quarter-bridge voltage and validity model.\\
\path{scripts/test_bridge_model_working.py} & Four offline bridge-model checks.\\
\path{soft_world/robot_description_working.py} & Six-disc, fifteen-joint RViz URDF.\\
\path{scripts/soft_limb_hang_working.py} & Hanging USD generator.\\
\path{scripts/soft_limb_working.py} & Position asset generator.\\
\path{scripts/create_six_disc_assets_working.py} & Native generator wrapper.\\
\path{scripts/update_six_disc_geometry_working.py} & Native authored geometry
update.\\
\path{usd/soft_limb_hanging_6_disc_working.usd} & Active hanging asset.\\
\path{usd/soft_limb_position_control_6_disc_working.usd} & Position asset.\\
\path{scripts/load_soft_limb_in_simulation_working.py} & Older two-limb backup.\\
\path{launch/soft_world_in_simulation_working.launch.py} & Backup launch.\\
\bottomrule
\end{longtable}

\section{Retained numerical excerpts}\label{app:results}
Working-controller test:
\begin{lstlisting}
VERIFY controller: base=36.0 deg,
split joints=[18.0, 18.0, 18.0, 18.0] deg;
tip=/World/SoftLimb/Disk_05
tip_orientation_error=0.0004 deg
last_second_phi_range=179.9996..179.9996 deg
intervals=[35.999966601826934, 35.99982317874544,
35.999893167570946, 35.999921731340976, 35.999989327436786]
VERIFY RESULT: PASS
\end{lstlisting}
Original twelve-disc controller with the working loader's adapter:
\begin{lstlisting}
Detected legacy 12-disc position controller:
rescaling bend targets by 11/5.
180 degrees now refers to the sixth/final disc Disk_05.
VERIFY controller: base=36.0 deg,
split joints=[18.0, 18.0, 18.0, 18.0] deg;
tip=/World/SoftLimb/Disk_05
tip_orientation_error=0.0004 deg
last_second_phi_range=179.9996..179.9996 deg
VERIFY RESULT: PASS
\end{lstlisting}
Separate ROS command verification:
\begin{lstlisting}
PASS actual ROS legacy controller:
base=36.000000, total=180.000000
\end{lstlisting}
The A and B arrays were both checked in the harness; the printed split array
above displays A only. Rounded equality of output fields is not a substitute
for a time-series dataset or uncertainty analysis.

\section{Suggested experiment-record template}
\begin{lstlisting}
Experiment ID / researcher / date:
Research question and hypothesis:
Git revision and uncommitted changes:
USD asset names, SHA-256 hashes, reference portability:
OS / ROS / Isaac / Python / GPU / driver versions:
Disc count, pitch, radius, masses:
Drive mode, stiffness, damping, body damping:
Timestep, solver iterations, gravity compensation:
Controller identity and parameters:
Loads, contact conditions, manual disturbances:
Simulation duration and timestamp convention:
Raw logs / rosbag / analysis script locations:
Metrics, tolerances, repeated-trial count:
Results, uncertainty, limitations, and disposition:
\end{lstlisting}
Record any generative-AI assistance in accordance with the researcher's
program and publication requirements. The researcher should review derivations,
citations, source changes, and claimed evidence before academic submission.

\end{document}
